\documentclass [12pt, a4paper] {extarticle}
\usepackage{setspace}
\usepackage[utf8x]{inputenc}
\usepackage[english,russian]{babel}
\usepackage[T2A]{fontenc}
\usepackage {graphicx}
\usepackage{ amssymb, latexsym, amsmath, amsthm}
\usepackage{geometry}
\geometry{left=3cm}
\geometry{right=1.5cm}
\geometry{top=2cm}
\geometry{bottom=3cm} 

\begin{document}

\thispagestyle {empty}

\begin {center}
    %\ \vspace{-3.5cm}

    \includegraphics [width = 0.5 \textwidth] {msu.png} \\
    {\scshape Московский Государственный Университет} \\
    Факультет вычислительной математики и кибернетики\\


    %\vfill
    %\vfill
    \vspace {2.7cm}
    {\LARGE Отчет по научно-исследовательской работе }

    %\vfill
    \vspace {0.5cm}
    {
        \LARGE \bfseries
        <<О сложности систем функций некоторых многозначных логик в классе
        поляризованных полиномов>> \\
    }

    %\vfill
    \vspace {0.7cm}
    {
        \Large
        студента 618/1 группы \\
        Гордеева Михаила Михайловича \\
    }

\end {center}


%\begin {center}
    %\Large
    %студента 618/1 группы \\
    %Гордеева Михаила Михайловича \\
%\end {center}

%\vfill
%\vfill
\vfill

%\begin {flushright}
    %\large
    %%Научный руководитель: <должность> Селезнева С.\,Н. \\
    %Научный руководитель: \\
    %доцент, к.ф.-м.н. Селезнева С.\,Н. \\
    %%Рекомендую  оценку -- <оценка> \underline{\hspace{138pt}} \\
    %Рекомендую  оценку -- 5 \underline{\hspace{50pt}} \\
    
    %\vspace {5mm}
%\end {flushright}

\hfill
\begin{tabular}{l}%{@{}l@{}}
    Научный руководитель: \\
    доцент, д.ф.-м.н. Селезнева С.\,Н. \\
    Рекомендую  оценку -- 5 \underline{\hspace{48pt}} \\
\end{tabular}

\vspace{1.5cm}

\begin {center}
Москва, 2016
\end {center}

%\vspace{1cm}

\enlargethispage {4 \baselineskip}

%\newpage
\setstretch{1.3}
\section{Введение}

Одним из стандартных способов задания функций k\nobreakdash-значной логики
являются поляризованные полиномиальные формы (ППФ), которые также называются
обобщенными формами Рида-Мюллера, или каноническими поляризованными полиномами.
В ППФ каждая переменная имеет определенную поляризацию.  Длиной полиномиальной
формы называется число попарно различных слагаемых в ней.  Длиной функции $f$ в
классе ППФ называется наименьшая длина среди длин всех поляризованных
полиномиальных форм, реализующих $f$. Функция Шеннона $L^K_k(n)$ длины
определяется как наибольшая длина среди всех функций $k$\nobreakdash значной
логики в классе $K$ от~$n$~переменных, если $K$ опущено, то подразумевается
класс ППФ. Практическое применение ППФ нашли при построении программируемых
логических матриц (ПЛМ)~\cite{ue04, sb90}, сложность ПЛМ напрямую зависит от
длины ППФ, по которой она построена. Поэтому в ряде работ исследуется сложность
ППФ различных функций \cite{sv93,pn95,ss02,kk05,sd08,mn12,sm09}. Также
рассматривают системы функций. Сложностью системы ППФ, имеющих один и тот же
вектор поляризации, называется число попарно различных слагаемых, встречающихся
во всех этих ППФ. Сложностью $L_k^\text{ППФ}(F)$ системы функций $k$-значной
логики $F = \{f_1(x_1, \ldots, x_n), \ldots, f_m(x_1, \ldots, x_n)\}$ в классе
ППФ называется минимальная сложность среди всех таких систем ППФ $\{p_1, \dots,
p_m\}$, что все ППФ $p_1, \dots, p_m$ имеют один и тот же вектор поляризации ,
и ППФ $p_j$ реализует функцию $f_j$, $j=1 \dots m$.  Понятно, что для
произвольной системы функций алгебры $k$-значной логики $F = \{f_1(x_1, \ldots,
x_n), \ldots, f_m(x_1, \ldots, x_n)\}$ верна оценка $L_k^\text{ППФ}(F)
\leqslant k^n$. Система фугнкций $F$ называется сложной, если
$L_k^\text{ППФ}(F) \geqslant k^n$. В данной работе найдены сложные системы
функций для некоторых $k > 3$.

В 1993  В.\,П.\,Супрун~\cite{sv93} получил первые оценки функции Шеннона для
функций алгебры логики :
$$
L_2(n) \geqslant C_n^{[\frac{n}{2}]},
$$
$$
L_2(n) < 3 \cdot 2^{n-1},
$$
где [$a$] обозначает целую часть $a$.

Точное значение функции Шеннона для функций алгебры логики в 1995\,г. было
найдено Н.\,А.\,Перязевым~\cite{pn95} :
$$
L_2(n) = \left[\frac{2^{n+1}}{3}\right].
$$

Функции $k$\nobreakdash-значных логик являются естественным обобщением функций
алгебры логики.  Для функций $k$\nobreakdash-значной логики верхняя оценка
функции Шеннона была получена в 2002\,г. С.\,Н.\,Селезневой~\cite{ss02} :
$$
L_k(n) < \frac{k(k-1)}{k(k-1)+1}k^n,
$$
в 2015 году она была улучшена \cite{by15}
$$
L_k(n) < \frac{k(k-1)-1}{k(k-1)}k^n.
$$


При построении ПЛМ рассматривают и другие полиномиальные формы. Например класс
обобщенных полиномиальных форм.  В классе обобщенных полиномиальных форм, в
отличие от класса поляризованных полиномиальных форм, переменные могут иметь
различную поляризацию в разных слагаемых. В статье
К.\,Д.\,Кириченко~\cite{kk05}, опубликованной в 2005\,г., был предложен метод
построения обобщенных полиномиальных форм из которого следует
$L^{\text{О.П.}}_2(n) = O(\frac{2 ^ n}{n})$.

В работах \cite{sd08, bs14} получено, что
$L^{\text{О.П.}}_k(n) = O(\frac{k ^ n}{n})$.

В 2012\,г. Н.\,К.\,Маркеловым была получена нижняя оценка функции Шеннона для
функции трехзначной логики в классе поляризованных полиномов~\cite{mn12}: 

\[ L_3(n) \geqslant \left[\frac{3}{4}3^n\right], \] 

в \cite{ss15} эта оценка была достигнута на симметрических функциях.  Также в
работе \cite{ss15} были получены некоторые оценки для систем функций.  А
именно, доказано существование таких систем $F_2$ и $F_3$ функций двузначной и
трехзначной логики соответственно, для которых

\[ L_2^\text{ППФ}(F_2) = 2^n \text{,} \] 

\[ L_3^\text{ППФ}(F_3) = 3^n \text{.} \]

\newpage
\section{Постановка задачи}
\begin{itemize}
    \item Найти сложные системы функций при некоторых $k$;

    \item Доказать существование сложных систем функций для все простых $k$.
\end{itemize}
    
\section{Выполненная работа}
\begin{itemize}
    \item Была написана программа для символьного преобразования полиномов;

    \item С помощью этой программы получены сложные системы функций для
        некоторых $k (5,7)$; 
        
    \item Также с помощью программы были проверены все системы функций
        определенного вида для все простых $k$ до 79;

    \item Доказано существование сложных систем для всех простых $k$.
\end{itemize}

\section{План работы на весенний семестр}
Завершить исследования и оформить выпускную работу.

\makeatletter
\renewcommand*{\@biblabel}[1]{\hfill#1.}
\makeatother

\newpage
\setstretch{1}

\begin{thebibliography}{0}
\bibitem{ue04} Угрюмов~Е.\,П. Цифровая схемотехника. СПб.: БХВ-Петербург, 2004.
\bibitem{sb90} Sasao T., Besslich P. On the complexity of mod-2 sum PLA’s  //
    IEEE Trans.on Comput. 39. N 2. 1990. P.~262--266.
\bibitem{sv93} Супрун~В.\,П. Сложность булевых функций в классе канонических
    поляризованных полиномов // Дискретная математика. 5.
    \textnumero 2. 1993. С. 111--115.
\bibitem{pn95} Перязев Н.\,А. Сложность булевых функций в классе полиномиальных
    поляризованных~форм // Алгебра и логика. 34.
    \textnumero 3. 1995. С. 323--326.
\bibitem{ss02} Селезнева С.\,H. О сложности представления функций многозначных
    логик поляризованными полиномами. Дискретная
    математика. 14. \textnumero 2. 2002. С.~48--53.
\bibitem{kk05} Кириченко~К.\,Д. Верхняя оценка сложности полиномиальных
    нормальных форм булевых функций
    // Дискретная математика. 17. \textnumero 3. 2005. С. 80--88.
\bibitem{sd08} Селезнева С.\,Н. Дайняк А.\,Б. О сложности обобщенных
    полиномов k\nobreakdash-значных функций // Вестник Московского
    университета. Серия 15. Вычислительная математика и кибернетика. 
        textnumero 3. 2008. С. 34--39.
\bibitem{mn12} Маркелов Н.\,К. Нижняя оценка сложности функций трехзначной
    логики в классе поляризованных полиномов // Вестник
    Московского университета. Серия 15. Вычислительная математика и
        кибернетика. \textnumero 3. 2012. С. 40--45.
\bibitem{sm09} Селезнева С.\,H. Маркелов Н.\,К. Быстрый алгоритм построения
    векторов коэффициэнтов поляризованных полиномов
    k-значных функций // Ученые записки Казанского университета. Серия
        Физико-математические науки. 2009. 151.
    \textnumero 2 С.~147-151.
\bibitem{bs14} Башов~М.\,А., Селезнева~С.\,Н. О длине функций $k$-значной
    логики в классе полиномиальных нормальных
    форм по модулю $k$ // Дискретная математика. -- 2014. -- Т. 26,
        вып. 3. -- С. 3-9.
\bibitem{ss15} Селезнева С.\,H. Сложность систем функций алгебры логики и
    систем функций трехзначной логики в классах
    поляризованных полиномиальных форм // Дискретная математика. -- 2015. --
        Т. 27, вып. 1. -- С. 111 -- 122. 
\bibitem{by15} Балюк~А.\,С., Янушковский~Г.\,В. Верхние оценки функций над
    конечными полями в некоторых классах
    кронекеровых форм // Известия Иркутского государственного университета.
        Серия: Математика. -- 2015. -- Т.14. -- С. 3-17.
\end{thebibliography}

\end{document}
